\section{Def rouleaux?}

Structure organisées de la couche limite planétaire. Peut etre mis en évidence par des rues de nuages écartées de $2-4 z_i$ et grosso modo dans le sens du vent (\cite{etlingRollVorticesPlanetary1993}). Les principales causes de ces structures 
\begin{itemize}
	\item \textbf{instabilité d'inflection}: lié au cisaillement de vent), 
	
	\item instabilité parallèles et des 
	
	\item \textbf{instabilité thermique}: convection de Rayleigh Bénard. Explique pk rolls surtout observés au dessus de la mer pendant CAO et au dessus de la terre durant la journée. 
	
\end{itemize}

\paragraph{Distinguer un rouleau}
paramètre d'autocroorélation spatiale R, puis facteur d'applatsissement f pour avoir l'allongement de la structure (\cite{lfarhImpactSurfaceTurbulent2024}).

\paragraph{Paramètre de stratification}
Paramètre de stabilité  $\zeta = -\frac{z_i}{L}$ : différencie les types de convection / influence flottabilité vs cisaillement.\\ 
$L$ est la longueur de Monin et Obukhov = - la distance où la turbulence produite par la flottabilité est supérieure à celle du cisaillement. \\
Pour $\zeta >> 1$ on est dans une situation purement convective. Des rouleaux sont oservés pour $\zeta \approx 0.5 < 1$ (\cite{lfarhImpactSurfaceTurbulent2024}) donc dans un BL neutre. \\
Les rouleaux sont attendus entre $6<\zeta<15$ (\cite{etlingRollVorticesPlanetary1993}), interaction entre convection et cisaillement. Mais aussi pour $\zeta \approx 2$ dans le cas d'une instabilité d'inflection. \\
C'est aussi souvent le cas pour $\zeta < 10$ ('free rolls') mais aussi pour $>20$ forced rolls
(\cite{gryschkaImpactForcedRoll2014}) 